% Maquetación común: mismos márgenes del documento y alineación consistente.
\usepackage{tabularx}
\newcolumntype{L}{>{\raggedright\arraybackslash}X}
\newcolumntype{C}{>{\centering\arraybackslash}X}
\DeclareCaptionFont{tabla}{\fontsize{11}{13}\selectfont}
\captionsetup[table]{font=tabla,position=top,skip=6pt,justification=raggedright,singlelinecheck=false}
\AtBeginEnvironment{table}{\centering\fontsize{11}{13}\selectfont\singlespacing\renewcommand{\arraystretch}{1.12}}
\AtEndEnvironment{table}{\par}
\setlength{\LTleft}{0pt plus 1fill}
\setlength{\LTright}{0pt plus 1fill}
\setlength{\LTpre}{8pt}
\setlength{\LTpost}{8pt}
% Los elementos flotantes pueden acompañar al texto sin quedarse en páginas solas.
\renewcommand{\topfraction}{0.90}
\renewcommand{\bottomfraction}{0.85}
\renewcommand{\textfraction}{0.08}
\renewcommand{\floatpagefraction}{0.80}
\setcounter{topnumber}{3}
\setcounter{bottomnumber}{2}
\setcounter{totalnumber}{4}
\setlength{\floatsep}{12pt plus 2pt minus 2pt}
\setlength{\textfloatsep}{14pt plus 2pt minus 2pt}
\setlength{\intextsep}{12pt plus 2pt minus 2pt}
% Páginas de flotantes alineadas arriba, sin huecos elásticos sobre una figura.
\makeatletter
\setlength{\@fptop}{0pt}
\setlength{\@fpsep}{14pt plus 2pt minus 2pt}
\setlength{\@fpbot}{0pt plus 1fil}
\makeatother
\AtBeginDocument{\raggedbottom\clubpenalty=10000\widowpenalty=10000\displaywidowpenalty=10000}
